% chapters/ch02e_logic_gates.tex
% Section 2.8: Logic Gates System

\section{Logic Gates System}
\label{sec:logic-gates-c2}

\nrmobox{Layer scope}{
This layer is a \emph{logical governance} system separate from Type
ZERO (which performs emotional shutdown / exit). It blocks assertion
errors caused by undefined definitions, undefined axes, set-theory
errors, or dismissed observations via a pre-output gate.
}

\subsection{The Four Core Gates}
\label{sec:dag-four-core-gates}

\begin{table}[ht]
\centering
\small
\begin{tabularx}{\linewidth}{lXXX}
\toprule
\textbf{Gate} & \textbf{Trigger condition} & \textbf{Pass condition} & \textbf{Action on fail} \\
\midrule
G1 Definition Gate
& Assertion with undefined concept
& Operational definition (measurable) is provided
& Hold judgement; request definition (or propose candidate) \\

G2 Axis Gate
& Mixed evaluation axes (speed / stability / creativity, etc.)
& Single axis, or if multi-axis: weights and order explicit
& Decompose axes; conclude per axis (no global assertion) \\

G3 Set Gate
& False dichotomy / part-to-whole fallacy
& Set relationships explicit; inference valid
& Halt inference; identify error type
(e.g., $\neg A \to B$, partial dominance $\to$ membership) \\

G4 Observation-First Gate
& General model overrides individual observation
& Only assumptions consistent with observation adopted
& Return to observation; recheck model premises \\
\bottomrule
\end{tabularx}
\caption[Logic Gates System]{Logic Gates System
(Table 2.3 --- Logical Governance Layer). The four gates form a
pipeline; an output passes only if all gates clear.}
\label{tab:logic-gates-c2}
\end{table}

\subsection{Ten-Line Statute (DAG Rules)}
\label{sec:dag-ten-line-statute}

For pasting into the SOP layer:

\begin{lstlisting}[basicstyle=\ttfamily\small,frame=single]
[DAG RULES]
1) Do not assign truth/falsity, superiority, or
   classification to undefined concepts (G1).
2) If evaluation axes are mixed, decompose axes and
   handle per single axis (G2).
3) Auto-block "not-A -> B" and "partial superiority ->
   whole membership" (G3).
4) When general model and observation collide,
   discard model and prioritize observation (G4).
5) Assertions include reversible conditions (premises);
   if premise unknown, defer.
6) Deferral is a mature output and is not weakness.
\end{lstlisting}

\subsection{Relationship with Type ZERO}
\label{sec:dag-vs-type-zero-c2}

\begin{table}[ht]
\centering
\small
\begin{tabularx}{\linewidth}{XX}
\toprule
\textbf{Type ZERO (emotional governance)}
& \textbf{DAG (logical governance)} \\
\midrule
Strong emotion / fatigue / impulse $\to$ stop, defer, retreat.
& Definition fluctuation / axis confusion / set error $\to$ block
assertion. \\
Purpose: avoid breakdown of decision-making
(psychological / behavioural side).
& Purpose: avoid breakdown of decision-making
(logical / argumentative side). \\
\bottomrule
\end{tabularx}
\caption[Type ZERO vs DAG]{Type ZERO and DAG operate on orthogonal
dimensions. They do not substitute for one another; both gates must
fire independently.}
\label{tab:type-zero-vs-dag-c2}
\end{table}

\subsection{Test Suite (Recurrence Prevention)}
\label{sec:dag-tests-c2}

\begin{description}
\item[Test A] ``Not a genius $\to$ therefore ordinary.''
\textit{Expected}: G3 stop ($\neg A \to B$ binary logic).
\item[Test B] ``Superior in partial domain $\to$ belongs to genius
domain.'' \textit{Expected}: G3 stop (partial-to-whole membership).
\item[Test C] Observation (high-output behaviour) and general model
(age-related decline) collide. \textit{Expected}: G4 discards model,
prioritises observation.
\end{description}

\subsection{Usage (Operational Notes)}
\label{sec:dag-usage-c2}

\nrmobox{Operational principle}{
The more the user wants ``assertion,'' the more DAG holds value. DAG
does not delay conclusions; it prevents broken conclusions from
becoming real.
}

\subsection{HST-N SOP --- Negative Assertion Pattern}
\label{sec:hstn-negative-assertion-c2}

The HST-N layer publishes an audit-only ``negative assertion''
record on a quarterly cadence. The pattern is structurally distinct
from a positive assertion (``everything is fine''): it states only
that under the current thresholds and observation windows,
\emph{no anomaly has been detected}.

\paragraph{Format.}

\begin{lstlisting}[basicstyle=\ttfamily\small,frame=single]
HST-N Negative Assertion (Quarterly)
- No anomaly detected under current thresholds.
- Observation windows: 4m / 12m (+ internal)
- Data availability: OK / Partial / Hold
\end{lstlisting}

\subsection{HST-N Observation Cycle}
\label{sec:hstn-observation-cycle-c2}

\paragraph{4-month state window (SOP-relevant).}
The 4-month window is the primary input to SOP-level Caution
detection.

\paragraph{12-month state window (Ruin-relevant).}
The 12-month window feeds Risk-Off detection (potential ruin
condition).

\paragraph{1-month and 6-month internal audit windows.}
The 1-month and 6-month windows are audit-only; they do not feed SOP
directly.

\subsection{Red Flag Codes}
\label{sec:hstn-red-flag-c2}

Red Flag codes are observational classifications, not actions. They
are forwarded to Type ZERO and SOP for downstream handling.

\begin{table}[ht]
\centering
\small
\begin{tabularx}{\linewidth}{llX}
\toprule
\textbf{Code} & \textbf{Name} & \textbf{Description} \\
\midrule
RF0 & Low-grade noise & Mild signal; logged but no action. \\
RF1 & Volatility spike & 4-month asset / benchmark volatility
elevated. \\
RF2 & Halt / Data hold & Shutdown signal; data not reliably
available. \\
RF3 & Liquidity gap & Bid-ask gap or resource gap. \\
RF4 & Correlation break & Cross-domain signal de-coupling
(observation only). \\
\bottomrule
\end{tabularx}
\caption[HST-N Red Flag codes]{HST-N Red Flag classification codes.
Codes are non-judgemental classifications passed to Type ZERO and
SOP for interpretation.}
\label{tab:hstn-red-flag-c2}
\end{table}

\paragraph{Log entry format.}

\begin{lstlisting}[basicstyle=\ttfamily\small,frame=single]
HST-N Red Flag Log
- Code:      RF1
- Window:    1m / 4m / 6m / 12m
- Asset:     (Ticker or Category)
- Timestamp: YYYY-MM-DD
- Note:      short text (no interpretation)
\end{lstlisting}

\subsection{Quarterly Frequency Summary}
\label{sec:hstn-frequency-summary-c2}

HST-N publishes a quarterly Red Flag frequency summary (audit only):

\begin{lstlisting}[basicstyle=\ttfamily\small,frame=single]
HST-N Frequency Summary (Quarterly)
- RF0: 5
- RF1: 1
- RF2: 0
- RF3: 0
- RF4: 1
\end{lstlisting}

\subsection{Hold Mode (Data Reliability)}
\label{sec:hstn-hold-c2}

When benchmark or asset data degrades, HST-N publishes a Hold
record:

\begin{itemize}
\item \texttt{OK} --- data reliable, full observation cycle
proceeds.
\item \texttt{Partial} --- partial coverage, classification flagged
partial.
\item \texttt{Hold} --- data unreliable, classification suspended,
SOP notified.
\end{itemize}

HST-N's ``Hold'' is a notification to SOP, not an executive action.
SOP retains decision authority.

\subsection{Update and Prohibitions}
\label{sec:hstn-update-prohibitions-c2}

\paragraph{Update conditions.}
The HST-N SOP is updated only when the upstream NRMO SOP changes.
Discretionary updates are prohibited.

\paragraph{Prohibitions.}

\begin{itemize}
\item State-transition execution: \textbf{prohibited}.
\item Direct action / proposal: \textbf{prohibited}.
\item TTM / PP execution: \textbf{prohibited}.
\item Issuing permitted-action sets: \textbf{prohibited}.
\end{itemize}

\paragraph{Design principle.}
``HST-N observes; HST-N does not decide.''
